\chapter*{راهنمای مطالعه}
\phantomsection
\addcontentsline{toc}{chapter}{راهنمای مطالعه}
\begingroup
\setlength{\parskip}{0.2em}
\ifCompleteNotes بخش اول\else این جزوه\fi{} برای یادگیری فصل ۱ و بخش‌های \secnum{۲.۱} تا \secnum{۲.۳} مرجع
\lr{Weste} و \lr{Harris} نوشته شده است. فصل صفر، ریاضی، مدار و منطق لازم را
از پایه آموزش می‌دهد. مفاهیم درس از مدل \lr{Switch} آغاز و سپس با جریان و ظرفیت
تکمیل می‌شوند. در هر مرحله، شرط‌های مدل پیش از محاسبه و پس از به‌دست‌آمدن پاسخ
بررسی می‌شوند \cite{weste,intro,mos}.

در مطالعه نخست، فصل صفر و سپس فصل‌های ۱ و ۲ را به ترتیب بخوانید.
برای هر مثال، اتصال‌ها، مرجع ولتاژ، واحدها و فرض ناحیه کاری را روی کاغذ بنویسید.
پرسش‌های هر درس پاسخ دارند. پس از تلاش مستقل، پاسخ را با حل تشریحی مقایسه کنید.
برای مرور، از جدول انتخاب مدل و پرسش‌های ترکیبی پایان جزوه استفاده کنید.

عنوان‌های اصلی فصل‌های ۱ و ۲ با بخش‌های کتاب هماهنگ‌اند.
زیربخش‌های آموزشی، به‌ویژه در منطق، ساخت و مثال پردازنده، برای یادگیری مرحله‌ای
گسترش یافته‌اند و شماره‌های داخلی جزوه دارند.
شماره روابط، شکل‌ها و جدول‌های جزوه از مرجع مستقل است.
ارجاع با عبارت «کتاب» یا «اسلاید» به شماره همان منبع اشاره می‌کند.
شماره‌گذاری تمرین‌های افزوده در هر مجموعه مستقل است. عنوان درس، مجموعه مربوط را مشخص می‌کند.

مثال‌های دارای داده فرضی برای توضیح روش تنظیم شده‌اند.
ابعاد و مدل‌های فناوری برگرفته از کتاب در زمینه همان مثال استفاده می‌شوند.
شکل‌های \lr{Layout} آموزشی‌اند و ارزیابی ساخت‌پذیری به قواعد فناوری و
بررسی‌های \lr{DRC} و \lr{LVS} نیاز دارد.
کدهای \lr{HDL} و \lr{SPICE} باید همراه با قرارداد، مدل و تحریک لازم خوانده شوند.

جدول راهنمای پوشش، محل موضوع‌های مرجع را در نسخه واحد نشان می‌دهد.
جدول اصلی و فهرست قابل کلیک، دو مسیر دسترسی به متن فراهم می‌کنند.
\begin{table}[ht]
\centering
\small
\textbf{جدول راهنمای پوشش موضوع‌های کتاب و اسلایدهای مرتبط}
\par\smallskip
\begin{tabularx}{\textwidth}{rXr}
\toprule
دامنه کتاب & موضوع و محل در جزوه & اسلاید مرتبط\\
\midrule
پیش‌نیاز افزوده & فصل \ref{chap:foundations}، محاسبه، مدار و منطق & مقدمات \cite{intro,mos}\\
\secnum{۱.۱} تا \secnum{۱.۳} & تاریخچه، مسیر طراحی و ساختار \lr{MOS} & \cite{intro} و آغاز \cite{mos}\\
\secnum{۱.۴} & بخش \ref{sec:book-cmos}، گیت، انتخاب و ذخیره داده & پایان \cite{intro}\\
\secnum{۱.۵} & بخش \ref{sec:book-fabrication}، ساخت و \lr{Layout} & ساختار در \cite{mos}\\
\secnum{۱.۶} & بخش \ref{sec:book-partitioning}، سطوح و دامنه‌ها & میانه \cite{intro}\\
\secnum{۱.۷} و \secnum{۱.۸} & بخش‌های \ref{sec:book-mips} و \ref{sec:book-logic-design}، پردازنده و \lr{HDL} & میانه \cite{intro}\\
\secnum{۱.۹} تا \secnum{۱.۱۲} & مدار، طراحی فیزیکی، بررسی و ساخت & تکمیل از کتاب\\
\secnum{۲.۱} & بخش \ref{sec:book-mos-structure}، سطح و کانال & \cite[اسلایدهای ۲ تا ۱۳]{mos}\\
\secnum{۲.۲} & بخش \ref{sec:book-mos-current}، استخراج جریان & \cite[اسلایدهای ۱۴ تا ۲۲]{mos}\\
\secnum{۲.۳} & بخش \ref{sec:book-cv}، ظرفیت و هندسه & \cite[اسلایدهای ۲۳ و ۲۴]{mos}\\
\bottomrule
\end{tabularx}
\end{table}

مبانی مدار فصل صفر با منابع تکمیلی معرفی‌شده در کتابنامه پشتیبانی شده‌اند
\cite{current,resistors,kirchhoff,capacitors}.
بخش \secnum{۲.۴} درباره اثرهای غیرایده‌آل و بخش \secnum{۲.۵} درباره مشخصه‌های
\lr{DC}، موضوع جزوه‌های دوم و سوم‌اند.
\endgroup
\clearpage
